\documentclass[12pt,a4paper]{article}
\usepackage[utf8]{vietnam}
\usepackage{amsmath,amssymb}
\usepackage{tabularx}
\usepackage{tikz}
\usepackage{xcolor}
\usepackage[top=0.6cm, bottom=0.6cm, left=0.6cm, right=0.6cm]{geometry}
\usepackage[dethi]{ex_test}

\SAOPTN{widthSA=2.8,heightSA=0.62}
\setlength{\parskip}{0pt}

% Dòng kẻ chấm mờ chuẩn viết tay
\newcommand{\dongke}[1]{%
  \begingroup%
  \setlength{\parskip}{0pt}%
  \baselineskip=18.5pt%
  \vspace{1.0mm}%
  \foreach \i in {1,...,#1}{%
    \noindent{\color{gray!60}\dotfill}\par%
  }%
  \vspace{0.6mm}%
  \endgroup%
}

% Định nghĩa ô bài tập tự luận trong bảng
\newcommand{\oBaiTap}[3][6]{%
  \vspace{1.0mm}%
  \noindent\textbf{#2} #3\par%
  \dongke{#1}%
}

\begin{document}
\pagestyle{empty}
% =============================================================================
% ========================= MA DE 101 - MAT TRUOC (TRANG 1/2) =============
% =============================================================================
\noindent
\begin{tabularx}{\linewidth}{|X|c|}
\hline
\textbf{PHIẾU KIỂM TRA ĐỊNH KỲ TOÁN 6 - TUẦN 10 - BUỔI 02} & \textbf{Mã đề: 101 - Trang 1/2} \\
\hline
\end{tabularx}

\vspace{1.0mm}

\noindent
\begin{tabularx}{\linewidth}{|p{3.2cm}|X|}
\hline
\centering\textbf{Điểm số} & \centering\textbf{Lời nhận xét của thầy cô} \tabularnewline
\rule{0pt}{1.8cm} & \\
\hline
\end{tabularx}

\vspace{1.0mm}

\noindent\textbf{Câu 1 (0,25 điểm).} Điền dấu thích hợp ($<, >, =$): $-15 \; \dots \; 0$ \shortans{$<$}\par\vspace{0.25mm}
\noindent\textbf{Câu 2 (0,25 điểm).} Điền dấu thích hợp ($<, >, =$): $-18 \; \dots \; -12$ \shortans{$<$}\par\vspace{0.25mm}
\noindent\textbf{Câu 3 (0,25 điểm).} Giá trị của biểu thức $|-19|$ là: \shortans{$19$}\par\vspace{0.25mm}
\noindent\textbf{Câu 4 (0,25 điểm).} Số nhỏ nhất trong dãy: $-12; 5; -8; 0; -15$ là: \shortans{$-15$}\par\vspace{0.25mm}
\noindent\textbf{Câu 5 (0,25 điểm).} Số nguyên liền trước của số $-8$ là: \shortans{$-9$}\par\vspace{0.25mm}
\noindent\textbf{Câu 6 (0,25 điểm).} Số nguyên âm lớn nhất có 2 chữ số là: \shortans{$-10$}\par\vspace{0.25mm}
\noindent\textbf{Câu 7 (0,25 điểm).} Có bao nhiêu số nguyên $x$ thỏa mãn $|x| = 7$? \shortans{$2$}\par\vspace{0.25mm}
\noindent\textbf{Câu 8 (0,25 điểm).} Hình thoi có tất cả bao nhiêu trục đối xứng? \shortans{$2$}\par\vspace{0.25mm}
\noindent\textbf{Câu 9 (0,25 điểm).} Hình thang cân có hai đáy $24\text{ cm}, 16\text{ cm}$, cao $10\text{ cm}$. Diện tích là: \shortans{$200\text{ cm}^2$}\par\vspace{0.25mm}
\noindent\textbf{Câu 10 (0,25 điểm).} Nhiệt độ Moscow $-8^\circ\text{C}$, Bắc Kinh $-3^\circ\text{C}$. Nơi nào lạnh hơn? \shortans{Moscow}\par\vspace{0.25mm}

\vspace{0.3mm}
\noindent\textbf{Câu 11 (3,0 điểm). Thực hiện phép tính (hợp lý nếu có thể):}
\vspace{0.4mm}

\noindent
\begin{tabularx}{\linewidth}{|X|X|}
\hline
\oBaiTap[7]{a)}{$48 + 157 + 52 + 43$} & \oBaiTap[7]{b)}{$38 \cdot 73 + 38 \cdot 27$} \\
\hline
\oBaiTap[7]{c)}{$5 \cdot 2^3 - |-27| : 3^2 + 2026^0$} & \oBaiTap[7]{d)}{$180 : (32 - 14) + 4^2 \cdot 3$} \\
\hline
\oBaiTap[7]{e)}{$200 : [45 - (12 - 7)^2]$} & \oBaiTap[7]{f)}{$160 : \{60 - [45 - (13 - 8)^2]\}$} \\
\hline
\end{tabularx}
\newpage
% =============================================================================
% ========================= MA DE 101 - MAT SAU (TRANG 2/2) ==============
% =============================================================================
\noindent
\begin{tabularx}{\linewidth}{|X|}
\hline
\centering\textbf{TOÁN 6 - TUẦN 10 - BUỔI 02 | Mã đề: 101 - Trang 2/2} \\
\hline
\end{tabularx}

\vspace{1.2mm}

\noindent\textbf{Câu 12 (2,0 điểm). Tìm số nguyên $x$, biết:}
\vspace{0.5mm}

\noindent
\begin{tabularx}{\linewidth}{|X|X|}
\hline
\oBaiTap[5]{a)}{$|x| = 8$} & \oBaiTap[5]{b)}{$|x - 3| = 5$} \\
\hline
\oBaiTap[7]{c)}{Liệt kê và tính tổng các số nguyên $x$: $-5 \le x < 6$} & \oBaiTap[7]{d)}{$(x + 5)(x - 3) = 0$} \\
\hline
\end{tabularx}

\vspace{1.2mm}

\noindent
\begin{tabularx}{\linewidth}{|X|}
\hline
\vspace{0.5mm}
\textbf{Câu 13 (2,0 điểm).} Bảng ghi nhận nhiệt độ lúc 6 giờ sáng mùa đông tại 5 thành phố lớn:\par\vspace{0.5mm}\centerline{\begin{tabular}{|l|c|c|c|c|c|}\hline \textbf{Thành phố} & Hà Nội & Bắc Kinh & Moscow & Tokyo & New York \\\hline \textbf{Nhiệt độ} & $16^\circ\text{C}$ & $-4^\circ\text{C}$ & $-12^\circ\text{C}$ & $5^\circ\text{C}$ & $-2^\circ\text{C}$ \\\hline\end{tabular}}\vspace{0.5mm}\par
\textbf{a) (1,0 điểm)} Sắp xếp nhiệt độ của 5 thành phố trên theo thứ tự từ thấp đến cao (tăng dần).\par
\textbf{b) (1,0 điểm)} Thành phố nào có nhiệt độ thấp nhất, thành phố nào cao nhất? Tính độ chênh lệch nhiệt độ giữa hai thành phố đó.\par
\vspace{0.4mm}
\noindent\textit{Bài giải:}\dongke{10}
\\ \hline
\end{tabularx}

\vspace{1.2mm}

\noindent
\begin{tabularx}{\linewidth}{|X|}
\hline
\vspace{0.5mm}
\textbf{Câu 14 (0,5 điểm).} Tìm tất cả các cặp số nguyên $(x, y)$ thỏa mãn phương trình: $(x - 1)(y + 2) = 5$.\par
\vspace{0.4mm}
\noindent\textit{Bài giải:}\dongke{7}
\\ \hline
\end{tabularx}

\newpage
% =============================================================================
% ========================= MA DE 102 - MAT TRUOC (TRANG 1/2) =============
% =============================================================================
\noindent
\begin{tabularx}{\linewidth}{|X|c|}
\hline
\textbf{PHIẾU KIỂM TRA ĐỊNH KỲ TOÁN 6 - TUẦN 10 - BUỔI 02} & \textbf{Mã đề: 102 - Trang 1/2} \\
\hline
\end{tabularx}

\vspace{1.0mm}

\noindent
\begin{tabularx}{\linewidth}{|p{3.2cm}|X|}
\hline
\centering\textbf{Điểm số} & \centering\textbf{Lời nhận xét của thầy cô} \tabularnewline
\rule{0pt}{1.8cm} & \\
\hline
\end{tabularx}

\vspace{1.0mm}

\noindent\textbf{Câu 1 (0,25 điểm).} Điền dấu thích hợp ($<, >, =$): $-24 \; \dots \; 0$ \shortans{$<$}\par\vspace{0.25mm}
\noindent\textbf{Câu 2 (0,25 điểm).} Điền dấu thích hợp ($<, >, =$): $-25 \; \dots \; -30$ \shortans{$>$}\par\vspace{0.25mm}
\noindent\textbf{Câu 3 (0,25 điểm).} Giá trị của biểu thức $|-27|$ là: \shortans{$27$}\par\vspace{0.25mm}
\noindent\textbf{Câu 4 (0,25 điểm).} Số nhỏ nhất trong dãy: $-9; 7; -14; 0; -6$ là: \shortans{$-14$}\par\vspace{0.25mm}
\noindent\textbf{Câu 5 (0,25 điểm).} Số nguyên liền sau của số $-15$ là: \shortans{$-14$}\par\vspace{0.25mm}
\noindent\textbf{Câu 6 (0,25 điểm).} Số nguyên âm nhỏ nhất có 2 chữ số là: \shortans{$-99$}\par\vspace{0.25mm}
\noindent\textbf{Câu 7 (0,25 điểm).} Có bao nhiêu số nguyên $x$ thỏa mãn $|x| = 9$? \shortans{$2$}\par\vspace{0.25mm}
\noindent\textbf{Câu 8 (0,25 điểm).} Tam giác đều có tất cả bao nhiêu trục đối xứng? \shortans{$3$}\par\vspace{0.25mm}
\noindent\textbf{Câu 9 (0,25 điểm).} Hình thang cân có hai đáy $25\text{ cm}, 15\text{ cm}$, cao $12\text{ cm}$. Diện tích là: \shortans{$240\text{ cm}^2$}\par\vspace{0.25mm}
\noindent\textbf{Câu 10 (0,25 điểm).} Nhiệt độ Seoul $-5^\circ\text{C}$, Tokyo $2^\circ\text{C}$. Nơi nào ấm hơn? \shortans{Tokyo}\par\vspace{0.25mm}

\vspace{0.3mm}
\noindent\textbf{Câu 11 (3,0 điểm). Thực hiện phép tính (hợp lý nếu có thể):}
\vspace{0.4mm}

\noindent
\begin{tabularx}{\linewidth}{|X|X|}
\hline
\oBaiTap[7]{a)}{$39 + 146 + 61 + 54$} & \oBaiTap[7]{b)}{$49 \cdot 65 + 49 \cdot 35$} \\
\hline
\oBaiTap[7]{c)}{$6 \cdot 2^3 - |-36| : 3^2 + 2025^0$} & \oBaiTap[7]{d)}{$160 : (29 - 13) + 5^2 \cdot 2$} \\
\hline
\oBaiTap[7]{e)}{$150 : [35 - (11 - 6)^2]$} & \oBaiTap[7]{f)}{$180 : \{60 - [55 - (14 - 9)^2]\}$} \\
\hline
\end{tabularx}
\newpage
% =============================================================================
% ========================= MA DE 102 - MAT SAU (TRANG 2/2) ==============
% =============================================================================
\noindent
\begin{tabularx}{\linewidth}{|X|}
\hline
\centering\textbf{TOÁN 6 - TUẦN 10 - BUỔI 02 | Mã đề: 102 - Trang 2/2} \\
\hline
\end{tabularx}

\vspace{1.2mm}

\noindent\textbf{Câu 12 (2,0 điểm). Tìm số nguyên $x$, biết:}
\vspace{0.5mm}

\noindent
\begin{tabularx}{\linewidth}{|X|X|}
\hline
\oBaiTap[5]{a)}{$|x| = 12$} & \oBaiTap[5]{b)}{$|x - 4| = 6$} \\
\hline
\oBaiTap[7]{c)}{Liệt kê và tính tổng các số nguyên $x$: $-6 < x \le 7$} & \oBaiTap[7]{d)}{$(x + 7)(x - 4) = 0$} \\
\hline
\end{tabularx}

\vspace{1.2mm}

\noindent
\begin{tabularx}{\linewidth}{|X|}
\hline
\vspace{0.5mm}
\textbf{Câu 13 (2,0 điểm).} Bảng ghi nhận nhiệt độ lúc 6 giờ sáng mùa đông tại 5 thành phố lớn:\par\vspace{0.5mm}\centerline{\begin{tabular}{|l|c|c|c|c|c|}\hline \textbf{Thành phố} & London & Paris & Berlin & Ottawa & Seoul \\\hline \textbf{Nhiệt độ} & $3^\circ\text{C}$ & $1^\circ\text{C}$ & $-5^\circ\text{C}$ & $-14^\circ\text{C}$ & $-3^\circ\text{C}$ \\\hline\end{tabular}}\vspace{0.5mm}\par
\textbf{a) (1,0 điểm)} Sắp xếp nhiệt độ của 5 thành phố trên theo thứ tự từ thấp đến cao (tăng dần).\par
\textbf{b) (1,0 điểm)} Thành phố nào có nhiệt độ thấp nhất, thành phố nào cao nhất? Tính độ chênh lệch nhiệt độ giữa hai thành phố đó.\par
\vspace{0.4mm}
\noindent\textit{Bài giải:}\dongke{10}
\\ \hline
\end{tabularx}

\vspace{1.2mm}

\noindent
\begin{tabularx}{\linewidth}{|X|}
\hline
\vspace{0.5mm}
\textbf{Câu 14 (0,5 điểm).} Tìm tất cả các cặp số nguyên $(x, y)$ thỏa mãn phương trình: $(x - 2)(y + 1) = 3$.\par
\vspace{0.4mm}
\noindent\textit{Bài giải:}\dongke{7}
\\ \hline
\end{tabularx}

\newpage
% =============================================================================
% ========================= MA DE 103 - MAT TRUOC (TRANG 1/2) =============
% =============================================================================
\noindent
\begin{tabularx}{\linewidth}{|X|c|}
\hline
\textbf{PHIẾU KIỂM TRA ĐỊNH KỲ TOÁN 6 - TUẦN 10 - BUỔI 02} & \textbf{Mã đề: 103 - Trang 1/2} \\
\hline
\end{tabularx}

\vspace{1.0mm}

\noindent
\begin{tabularx}{\linewidth}{|p{3.2cm}|X|}
\hline
\centering\textbf{Điểm số} & \centering\textbf{Lời nhận xét của thầy cô} \tabularnewline
\rule{0pt}{1.8cm} & \\
\hline
\end{tabularx}

\vspace{1.0mm}

\noindent\textbf{Câu 1 (0,25 điểm).} Điền dấu thích hợp ($<, >, =$): $0 \; \dots \; -35$ \shortans{$>$}\par\vspace{0.25mm}
\noindent\textbf{Câu 2 (0,25 điểm).} Điền dấu thích hợp ($<, >, =$): $-42 \; \dots \; -37$ \shortans{$<$}\par\vspace{0.25mm}
\noindent\textbf{Câu 3 (0,25 điểm).} Giá trị của biểu thức $|-35|$ là: \shortans{$35$}\par\vspace{0.25mm}
\noindent\textbf{Câu 4 (0,25 điểm).} Số lớn nhất trong dãy: $-16; -3; -20; 0; -1$ là: \shortans{$0$}\par\vspace{0.25mm}
\noindent\textbf{Câu 5 (0,25 điểm).} Số nguyên liền trước của số $-12$ là: \shortans{$-13$}\par\vspace{0.25mm}
\noindent\textbf{Câu 6 (0,25 điểm).} Số nguyên âm lớn nhất là số nào? \shortans{$-1$}\par\vspace{0.25mm}
\noindent\textbf{Câu 7 (0,25 điểm).} Số nguyên $x$ thỏa mãn $|x| = 0$ là: \shortans{$0$}\par\vspace{0.25mm}
\noindent\textbf{Câu 8 (0,25 điểm).} Hình tròn có tất cả bao nhiêu trục đối xứng? \shortans{Vô số}\par\vspace{0.25mm}
\noindent\textbf{Câu 9 (0,25 điểm).} Hình thang cân có hai đáy $28\text{ cm}, 12\text{ cm}$, cao $15\text{ cm}$. Diện tích là: \shortans{$300\text{ cm}^2$}\par\vspace{0.25mm}
\noindent\textbf{Câu 10 (0,25 điểm).} Ngăn đông $-18^\circ\text{C}$, ngăn mát $4^\circ\text{C}$. Ngăn nào lạnh hơn? \shortans{Ngăn đông}\par\vspace{0.25mm}

\vspace{0.3mm}
\noindent\textbf{Câu 11 (3,0 điểm). Thực hiện phép tính (hợp lý nếu có thể):}
\vspace{0.4mm}

\noindent
\begin{tabularx}{\linewidth}{|X|X|}
\hline
\oBaiTap[7]{a)}{$53 + 138 + 47 + 62$} & \oBaiTap[7]{b)}{$58 \cdot 82 + 58 \cdot 18$} \\
\hline
\oBaiTap[7]{c)}{$7 \cdot 2^3 - |-45| : 3^2 + 2024^0$} & \oBaiTap[7]{d)}{$210 : (36 - 15) + 3^2 \cdot 6$} \\
\hline
\oBaiTap[7]{e)}{$280 : [65 - (13 - 8)^2]$} & \oBaiTap[7]{f)}{$200 : \{70 - [55 - (12 - 7)^2]\}$} \\
\hline
\end{tabularx}
\newpage
% =============================================================================
% ========================= MA DE 103 - MAT SAU (TRANG 2/2) ==============
% =============================================================================
\noindent
\begin{tabularx}{\linewidth}{|X|}
\hline
\centering\textbf{TOÁN 6 - TUẦN 10 - BUỔI 02 | Mã đề: 103 - Trang 2/2} \\
\hline
\end{tabularx}

\vspace{1.2mm}

\noindent\textbf{Câu 12 (2,0 điểm). Tìm số nguyên $x$, biết:}
\vspace{0.5mm}

\noindent
\begin{tabularx}{\linewidth}{|X|X|}
\hline
\oBaiTap[5]{a)}{$|x| = 15$} & \oBaiTap[5]{b)}{$|x - 2| = 7$} \\
\hline
\oBaiTap[7]{c)}{Liệt kê và tính tổng các số nguyên $x$: $-4 \le x \le 5$} & \oBaiTap[7]{d)}{$(x + 6)(x - 2) = 0$} \\
\hline
\end{tabularx}

\vspace{1.2mm}

\noindent
\begin{tabularx}{\linewidth}{|X|}
\hline
\vspace{0.5mm}
\textbf{Câu 13 (2,0 điểm).} Bảng ghi nhận nhiệt độ lúc 6 giờ sáng mùa đông tại 5 thành phố lớn:\par\vspace{0.5mm}\centerline{\begin{tabular}{|l|c|c|c|c|c|}\hline \textbf{Thành phố} & Madrid & Warsaw & Helsinki & Kyiv & Rome \\\hline \textbf{Nhiệt độ} & $8^\circ\text{C}$ & $-6^\circ\text{C}$ & $-11^\circ\text{C}$ & $-7^\circ\text{C}$ & $10^\circ\text{C}$ \\\hline\end{tabular}}\vspace{0.5mm}\par
\textbf{a) (1,0 điểm)} Sắp xếp nhiệt độ của 5 thành phố trên theo thứ tự từ thấp đến cao (tăng dần).\par
\textbf{b) (1,0 điểm)} Thành phố nào có nhiệt độ thấp nhất, thành phố nào cao nhất? Tính độ chênh lệch nhiệt độ giữa hai thành phố đó.\par
\vspace{0.4mm}
\noindent\textit{Bài giải:}\dongke{10}
\\ \hline
\end{tabularx}

\vspace{1.2mm}

\noindent
\begin{tabularx}{\linewidth}{|X|}
\hline
\vspace{0.5mm}
\textbf{Câu 14 (0,5 điểm).} Tìm tất cả các cặp số nguyên $(x, y)$ thỏa mãn phương trình: $(x + 1)(y - 3) = 7$.\par
\vspace{0.4mm}
\noindent\textit{Bài giải:}\dongke{7}
\\ \hline
\end{tabularx}

\newpage
% =============================================================================
% ========================= MA DE 104 - MAT TRUOC (TRANG 1/2) =============
% =============================================================================
\noindent
\begin{tabularx}{\linewidth}{|X|c|}
\hline
\textbf{PHIẾU KIỂM TRA ĐỊNH KỲ TOÁN 6 - TUẦN 10 - BUỔI 02} & \textbf{Mã đề: 104 - Trang 1/2} \\
\hline
\end{tabularx}

\vspace{1.0mm}

\noindent
\begin{tabularx}{\linewidth}{|p{3.2cm}|X|}
\hline
\centering\textbf{Điểm số} & \centering\textbf{Lời nhận xét của thầy cô} \tabularnewline
\rule{0pt}{1.8cm} & \\
\hline
\end{tabularx}

\vspace{1.0mm}

\noindent\textbf{Câu 1 (0,25 điểm).} Điền dấu thích hợp ($<, >, =$): $0 \; \dots \; -48$ \shortans{$>$}\par\vspace{0.25mm}
\noindent\textbf{Câu 2 (0,25 điểm).} Điền dấu thích hợp ($<, >, =$): $-56 \; \dots \; -65$ \shortans{$>$}\par\vspace{0.25mm}
\noindent\textbf{Câu 3 (0,25 điểm).} Giá trị của biểu thức $|-48|$ là: \shortans{$48$}\par\vspace{0.25mm}
\noindent\textbf{Câu 4 (0,25 điểm).} Số lớn nhất trong dãy: $-25; -7; -30; -2; -11$ là: \shortans{$-2$}\par\vspace{0.25mm}
\noindent\textbf{Câu 5 (0,25 điểm).} Số nguyên liền sau của số $-20$ là: \shortans{$-19$}\par\vspace{0.25mm}
\noindent\textbf{Câu 6 (0,25 điểm).} Số nguyên âm nhỏ nhất có 1 chữ số là: \shortans{$-9$}\par\vspace{0.25mm}
\noindent\textbf{Câu 7 (0,25 điểm).} Có bao nhiêu số nguyên $x$ thỏa mãn $|x| = -5$? \shortans{$0$}\par\vspace{0.25mm}
\noindent\textbf{Câu 8 (0,25 điểm).} Hình lục giác đều có tất cả bao nhiêu trục đối xứng? \shortans{$6$}\par\vspace{0.25mm}
\noindent\textbf{Câu 9 (0,25 điểm).} Hình thang cân có hai đáy $30\text{ cm}, 14\text{ cm}$, cao $10\text{ cm}$. Diện tích là: \shortans{$220\text{ cm}^2$}\par\vspace{0.25mm}
\noindent\textbf{Câu 10 (0,25 điểm).} Sa Pa $-1^\circ\text{C}$, Mẫu Sơn $-4^\circ\text{C}$. Nơi nào rét hơn? \shortans{Mẫu Sơn}\par\vspace{0.25mm}

\vspace{0.3mm}
\noindent\textbf{Câu 11 (3,0 điểm). Thực hiện phép tính (hợp lý nếu có thể):}
\vspace{0.4mm}

\noindent
\begin{tabularx}{\linewidth}{|X|X|}
\hline
\oBaiTap[7]{a)}{$26 + 165 + 74 + 35$} & \oBaiTap[7]{b)}{$67 \cdot 76 + 67 \cdot 24$} \\
\hline
\oBaiTap[7]{c)}{$8 \cdot 2^3 - |-18| : 3^2 + 2027^0$} & \oBaiTap[7]{d)}{$240 : (42 - 18) + 4^2 \cdot 4$} \\
\hline
\oBaiTap[7]{e)}{$300 : [75 - (14 - 9)^2]$} & \oBaiTap[7]{f)}{$240 : \{80 - [65 - (15 - 10)^2]\}$} \\
\hline
\end{tabularx}
\newpage
% =============================================================================
% ========================= MA DE 104 - MAT SAU (TRANG 2/2) ==============
% =============================================================================
\noindent
\begin{tabularx}{\linewidth}{|X|}
\hline
\centering\textbf{TOÁN 6 - TUẦN 10 - BUỔI 02 | Mã đề: 104 - Trang 2/2} \\
\hline
\end{tabularx}

\vspace{1.2mm}

\noindent\textbf{Câu 12 (2,0 điểm). Tìm số nguyên $x$, biết:}
\vspace{0.5mm}

\noindent
\begin{tabularx}{\linewidth}{|X|X|}
\hline
\oBaiTap[5]{a)}{$|x| = 9$} & \oBaiTap[5]{b)}{$|x - 5| = 4$} \\
\hline
\oBaiTap[7]{c)}{Liệt kê và tính tổng các số nguyên $x$: $-7 < x < 6$} & \oBaiTap[7]{d)}{$(x + 8)(x - 5) = 0$} \\
\hline
\end{tabularx}

\vspace{1.2mm}

\noindent
\begin{tabularx}{\linewidth}{|X|}
\hline
\vspace{0.5mm}
\textbf{Câu 13 (2,0 điểm).} Bảng ghi nhận nhiệt độ lúc 6 giờ sáng mùa đông tại 5 thành phố lớn:\par\vspace{0.5mm}\centerline{\begin{tabular}{|l|c|c|c|c|c|}\hline \textbf{Thành phố} & Chicago & Toronto & Oslo & Vienna & Bangkok \\\hline \textbf{Nhiệt độ} & $-9^\circ\text{C}$ & $-8^\circ\text{C}$ & $-13^\circ\text{C}$ & $-1^\circ\text{C}$ & $25^\circ\text{C}$ \\\hline\end{tabular}}\vspace{0.5mm}\par
\textbf{a) (1,0 điểm)} Sắp xếp nhiệt độ của 5 thành phố trên theo thứ tự từ thấp đến cao (tăng dần).\par
\textbf{b) (1,0 điểm)} Thành phố nào có nhiệt độ thấp nhất, thành phố nào cao nhất? Tính độ chênh lệch nhiệt độ giữa hai thành phố đó.\par
\vspace{0.4mm}
\noindent\textit{Bài giải:}\dongke{10}
\\ \hline
\end{tabularx}

\vspace{1.2mm}

\noindent
\begin{tabularx}{\linewidth}{|X|}
\hline
\vspace{0.5mm}
\textbf{Câu 14 (0,5 điểm).} Tìm tất cả các cặp số nguyên $(x, y)$ thỏa mãn phương trình: $(x + 2)(y - 1) = 11$.\par
\vspace{0.4mm}
\noindent\textit{Bài giải:}\dongke{7}
\\ \hline
\end{tabularx}
\end{document}
